\documentclass{ucshinforme}

% --- Datos de la portada y del encabezado ---
\asignatura{Bases de Datos}
\anio{2026}
\semestre{2}
\titulo{Título del informe}
\autores{Nombre Apellido\\ Segundo Integrante\\ Tercer Integrante}
\profesor{Nombre del Profesor}
\facultad{Facultad de Ingeniería}
\carrera{Ingeniería Civil en Computación e Informática}
% \fechaentrega{15 de octubre de 2026}

% --- Archivo con las fuentes ---
\addbibresource{referencias.bib}

\begin{document}

\makeportada

% --- Índices (singlespace: se ven más compactos que con 1.5) ---
\begin{singlespace}
\tableofcontents
\clearpage
\listoffigures    % borra esta línea y su \clearpage si no usas figuras
\clearpage
\listoftables     % borra esta línea y su \clearpage si no usas tablas
\clearpage
\end{singlespace}

% --- Contenido ---
\section{Introducción}
Texto de ejemplo, con una cita \parencite{elmasri2016}.

\section{Desarrollo}
\subsection{Marco teórico}
Texto de ejemplo, con otra cita \parencite{date2019}.

\subsection{Resultados}
La Figura~\ref{fig:ejemplo} y la Tabla~\ref{tab:ejemplo} son ejemplos.

\begin{figure}[h]
  \centering
  \fbox{\parbox{6cm}{\centering\vspace{1.5cm}Imagen de ejemplo\vspace{1.5cm}}}
  \caption{Ejemplo de figura.}
  \label{fig:ejemplo}
\end{figure}

\begin{table}[h]
  \centering
  \begin{tabular}{lc}
    \hline
    Concepto & Valor \\
    \hline
    A & 10 \\
    B & 20 \\
    \hline
  \end{tabular}
  \caption{Ejemplo de tabla.}
  \label{tab:ejemplo}
\end{table}

\section{Conclusiones}
Texto de ejemplo...

\newpage
\printbibliography[heading=bibintoc,title=Referencias]

\end{document}